\documentclass[12pt,a4paper,titlepage]{article}

\usepackage{amsmath}
\usepackage{titlesec}
\usepackage{graphicx}
\usepackage[portuguese]{babel}
\usepackage[colorlinks=true,allcolors=black]{hyperref}

\graphicspath{ {./../../../../} }

\newcommand{\sectionbreak}{\clearpage}
\renewcommand{\arraystretch}{1.5}

\title{Transformada de Laplace Inversa}
\author{
    \includegraphics[scale=0.9]{puc}
    \\[1.5em]
    \\Fernando de Lima Gonçalves
}
\date{Belo Horizonte\\2026}

\begin{document}
\maketitle
\tableofcontents

\section{Revisão}
\subsection{Transformadas de Laplace}
\[
    \boxed{L\{f\relax(t)\} = \int_{0}^\infty e^{-st}\,f\relax(t)\,dt \;,\; t \geq 0}
\]

\begin{center}
    \begin{tabular}{||c|c||}
        \hline
        $f\relax(t)$ & $F\relax(s) = L\{f\relax(t)\}$\\
        \hline\hline
        1 & $\frac1 s$\\
        \hline
        $t^{n}$ & $\frac{n!}{s^{n+1}}$\\
        \hline
        $e^{a\,t}$ & $\frac{1}{s-a}$\\
        \hline
        $\cos(at)$ & $\frac{s}{s^{2}+a^{2}}$\\
        \hline
        $\sin(at)$ & $\frac{a}{s^{2}+a^{2}}$\\
        \hline
        $t^{n}\,e^{at}$ & $\frac{n!}{(s-a)^{n+1}}$\\
        \hline\hline
    \end{tabular}
\end{center}

\subsection{Frações Parciais}
Definimos \textit{Frações Parciais} como uma resolução de integrais por meio de:

\section{Exercícios}

\section{Imagens}

\section{Referências}

\end{document}
